\usepackage[T1]{fontenc}
\usepackage[utf8]{inputenc}
\usepackage{mathpazo}
% Typewriter text in the txfonts typewriter, a vector font (the default EC one is a bitmap without cm-super),
% so that every font is an embedded Type 1 font, as print shops ask.
\renewcommand{\ttdefault}{txtt}
% Embedded vector sans-serif labels, consistent with the covers.
\renewcommand{\sfdefault}{qhv}
\usepackage{amsmath,amssymb,amsthm,mathtools}
% Every PDF uses the standard 7 x 10 inch book layout. Its identical file also
% serves as the paperback interior; the full cover is a separate publishing asset.
% Generated by tools/print_layout.py; edit print-layout.json.
\newcommand{\PRINTWIDTH}{7}
\newcommand{\PRINTHEIGHT}{10}
\newcommand{\PRINTBLEED}{0.125}
\newcommand{\PRINTSPINECOEFF}{0.002347}

\providecommand{\PRINTINNER}{0.75in}
% Keep these hooks for volume layouts and stacked mathematical figures.
\newif\ifprintformat\printformattrue
\newif\ifsmallformat\smallformattrue
\makeatletter\@twosidetrue\@mparswitchtrue\@openrighttrue\makeatother
\usepackage[paperwidth=\PRINTWIDTH in,paperheight=\PRINTHEIGHT in,twoside,inner=\PRINTINNER,outer=0.6in,
  top=0.75in,bottom=0.8in,headsep=0.2in,footskip=0.4in]{geometry}
% Break long formula chains for the book's text measure and constrain displays.
\newcommand{\qq}{\\}\newcommand{\eqbr}{\\&}
\AtBeginDocument{%
  \def\[#1\]{\begin{equation*}\sbox0{$\displaystyle#1$}%
    \ifdim\wd0>\linewidth\typeout{SHRINK \the\wd0\space > \the\linewidth\space at line \the\inputlineno}%
      \resizebox{\linewidth}{!}{\usebox0}\else\usebox0\fi\end{equation*}}}
\usepackage[final]{microtype}
\setlength{\emergencystretch}{3em}
\usepackage{enumitem}
\usepackage[dvipsnames]{xcolor}
% Approved charcoal/coral editorial identity. See docs/DESIGN.md.
\definecolor{ocBlack}{HTML}{191919}
\definecolor{ocCoral}{HTML}{F28B79}
\definecolor{ocCoralDeep}{HTML}{A34232}
\definecolor{ocInk}{HTML}{252323}
\definecolor{ocCream}{HTML}{F5F1E9}
% Compatibility names for mathematical figures authored before the palette change.
\colorlet{ocGold}{ocCoral}
\colorlet{ocGoldDeep}{ocCoralDeep}
% Legacy names remain aliases so semantic boxes and existing figures stay compatible.
\colorlet{ocTerracotta}{ocCoralDeep}
\colorlet{ocRust}{ocCoralDeep}
\definecolor{ocSepia}{HTML}{795B54}
\definecolor{ocOlive}{HTML}{5D6D65}
\definecolor{ocUmber}{HTML}{746072}
\colorlet{ocExercise}{ocSepia}
\colorlet{ocReminder}{ocCoralDeep}
\colorlet{ocBeauty}{ocOlive}
\colorlet{ocThink}{ocUmber}

% Paperback interiors are black ink on white paper; tint hierarchy survives in gray.
\ifprintformat
  \colorlet{ocExercise}{black!65}
  \colorlet{ocReminder}{black!70}
  \colorlet{ocBeauty}{black!60}
  \colorlet{ocThink}{black!75}
\fi
\usepackage[most]{tcolorbox}
\usepackage{needspace}
\usepackage{currfile}
% Book pages keep framed material whole. Record natural heights for the print
% audit; oversized material must be reorganised in the source, never clipped.
\makeatletter
\newwrite\ocLayoutFile
\AtBeginDocument{\immediate\openout\ocLayoutFile=\jobname.oclayout\relax}
\AtEndDocument{\immediate\closeout\ocLayoutFile}
\newcommand{\ocCheckBox}{%
  \immediate\write\ocLayoutFile{BOX|\currfilepath|\the\inputlineno|\tcb@height|\the\textheight}%
  \ifdim\tcb@height>\textheight
    \PackageError{orderchance}{Box exceeds the 7x10 text height at \currfilepath:\the\inputlineno}%
      {Reorganise this material into shorter, logically complete boxes. Do not shrink the type.}%
  \fi}
\newcommand{\ocSplitBoxError}{%
  \PackageError{orderchance}{A framed box was split across pages}%
    {Use an unbreakable box and reorganise oversized source material.}}
\tcbset{finish unbroken={\ocCheckBox},
  finish first={\ocSplitBoxError},finish middle={\ocSplitBoxError},finish last={\ocSplitBoxError}}
% Needspace inserts a break opportunity. Preserve LaTeX's heading keep when the
% first hint or solution immediately follows a section heading.
\newcommand{\ocOpeningSpace}{\if@nobreak\else\Needspace{3\baselineskip}\fi}
\makeatother
% Avoid isolated opening/closing lines and an isolated line before a display.
\clubpenalty=10000
\widowpenalty=10000
\displaywidowpenalty=10000
\raggedbottom
% Packages and macros used by figures; shared by the books and tools/figure.tex.
\usepackage{tikz}
\usetikzlibrary{arrows.meta,calc,positioning,decorations.pathreplacing}
\usepackage{pgfplots}
\pgfplotsset{compat=1.18}

\usepackage{pdfpages}
\usepackage{adjustbox}
% The print interior starts with its half-title. The build attaches the approved
% front to a separate reading PDF; EPUB and paperback wraps have their own covers.
% Retain the manuscript hook so every main.tex keeps the same front-matter order.
\newcommand{\makecover}{}
% Running heads in print: the chapter title on left pages, the section on right pages, in small
% italics, cut short with an ellipsis when too long for the text measure.
\ifprintformat
  \usepackage{fancyhdr}
  \usepackage[fit]{truncate}
  \pagestyle{fancy}\fancyhf{}
  \renewcommand{\headrulewidth}{0pt}
  \renewcommand{\chaptermark}[1]{\markboth{#1}{#1}}
  \renewcommand{\sectionmark}[1]{\markright{\thesection\quad #1}}
  \fancyhead[LE]{\small\thepage\hfill\truncate{0.85\headwidth}{\itshape\leftmark}}
  \fancyhead[RO]{\small\truncate{0.85\headwidth}{\itshape\rightmark}\hfill\thepage}
  % The blank page before a chapter has no running head.
  \makeatletter
  \renewcommand{\cleardoublepage}{\clearpage\if@twoside\ifodd\c@page\else
    \hbox{}\thispagestyle{empty}\newpage\fi\fi}
  \makeatother
\fi
% PDF links are black; online and EPUB themes are configured separately.
\ifprintformat
  \usepackage[hidelinks]{hyperref}
\else
  \usepackage[colorlinks=true,linkcolor=ocRust,urlcolor=ocRust]{hyperref}
\fi
% Half-title, series page, title page and copyright page (\maketitle builds all four).
% Front matter of every volume, in the order of a published book. \maketitle (called by each
% main.tex) builds four pages: the half-title (recto), the series page (verso), the title page
% (recto) and the copyright page (verso). The preface and the table of contents follow.
%
% Series-wide settings (one place for every volume):
\newcommand{\seriesname}{Order and Chance}
\newcommand{\seriessubtitle}{The Hidden Architecture of Mathematics}
% Generated from volume-subtitles.json by tools/series_metadata.py.
\newcommand{\volumesubtitledef}[2]{\expandafter\def\csname ocsubtitle@#1\endcsname{#2}}
\newcommand{\volumesubtitle}[1]{\ifcsname ocsubtitle@#1\endcsname\csname ocsubtitle@#1\endcsname\fi}
\volumesubtitledef{0}{Proof, Problem Solving and Mathematical Thinking}
\volumesubtitledef{1}{Primes, Patterns and Proofs}
\volumesubtitledef{2}{The Language of Symmetry}
\volumesubtitledef{3}{The Logic of Uncertainty}
\volumesubtitledef{4}{From Formula to Insight}
\volumesubtitledef{5}{Where Numbers Meet Infinity}
\volumesubtitledef{6}{A Landscape of Limits}
\volumesubtitledef{7}{Limits, Measures and Transformations}
\volumesubtitledef{8}{The Geometry of Connection}
\volumesubtitledef{9}{Categories, Rings and Homology}
\volumesubtitledef{10}{The Structure of Randomness}
\volumesubtitledef{11}{The Geometry of Random Paths}
\volumesubtitledef{12}{From Equations to Hidden Worlds}
\volumesubtitledef{13}{The Hidden Rhythm of Primes}
\volumesubtitledef{14}{Where Number Meets Randomness}
\volumesubtitledef{P1}{Problems Across Mathematical Borders}
\newcommand{\volumecoversubtitle}[1]{\csname occoversubtitle@#1\endcsname}
\expandafter\def\csname occoversubtitle@0\endcsname{Proof, Problem Solving\\and Mathematical Thinking}
\expandafter\def\csname occoversubtitle@1\endcsname{Primes, Patterns and Proofs}
\expandafter\def\csname occoversubtitle@2\endcsname{The Language of Symmetry}
\expandafter\def\csname occoversubtitle@3\endcsname{The Logic of Uncertainty}
\expandafter\def\csname occoversubtitle@4\endcsname{From Formula to Insight}
\expandafter\def\csname occoversubtitle@5\endcsname{Where Numbers Meet Infinity}
\expandafter\def\csname occoversubtitle@6\endcsname{A Landscape of Limits}
\expandafter\def\csname occoversubtitle@7\endcsname{Limits, Measures and Transformations}
\expandafter\def\csname occoversubtitle@8\endcsname{The Geometry of Connection}
\expandafter\def\csname occoversubtitle@9\endcsname{Categories, Rings and Homology}
\expandafter\def\csname occoversubtitle@10\endcsname{The Structure of Randomness}
\expandafter\def\csname occoversubtitle@11\endcsname{The Geometry of Random Paths}
\expandafter\def\csname occoversubtitle@12\endcsname{From Equations to Hidden Worlds}
\expandafter\def\csname occoversubtitle@13\endcsname{The Hidden Rhythm of Primes}
\expandafter\def\csname occoversubtitle@14\endcsname{Where Number Meets Randomness}
\expandafter\def\csname occoversubtitle@P1\endcsname{Problems Across Mathematical Borders}

% Generated from common/cover-notes.json by tools/series_cover.py.
\expandafter\def\csname occovercredit@1\endcsname{Cover illustration: the Collatz tree from 2 up to 8,000; see About the Cover.}
\expandafter\def\csname occovercredit@2\endcsname{Cover illustration: the Ulam spiral of the primes up to the square of 71; see About the Cover.}
\expandafter\def\csname occovercredit@3\endcsname{Cover illustration: the icosahedron and its axes of symmetry; see About the Cover.}
\expandafter\def\csname occovercredit@4\endcsname{Cover illustration: Galton's quincunx and the binomial bell; see About the Cover.}
\expandafter\def\csname occovercredit@5\endcsname{Cover illustration: the Euler spiral, also called the Cornu spiral, with its osculating circles; see About the Cover.}
\expandafter\def\csname occovercredit@6\endcsname{Cover illustration: Ford circles and the Farey tessellation; see About the Cover.}
\expandafter\def\csname occovercredit@7\endcsname{Cover illustration: the Cantor function over the Cantor set; see About the Cover.}
\expandafter\def\csname occovercredit@8\endcsname{Cover illustration: Fourier sums of the square wave, and its spectrum; see About the Cover.}
\expandafter\def\csname occovercredit@9\endcsname{Cover illustration: the hyperbolic tiling by heptagons, three at a vertex; see About the Cover.}
\expandafter\def\csname occovercredit@10\endcsname{Cover illustration: the snake lemma and its connecting homomorphism; see About the Cover.}
\expandafter\def\csname occovercredit@11\endcsname{Cover illustration: simple random walks and their diffusive envelope; see About the Cover.}
\expandafter\def\csname occovercredit@12\endcsname{Cover illustration: a Brownian path stopped on the unit circle; see About the Cover.}
\expandafter\def\csname occovercredit@13\endcsname{Cover illustration: an elliptic curve and the chord-and-tangent law; see About the Cover.}
\expandafter\def\csname occovercredit@14\endcsname{Cover illustration: the zeta function walking the critical line; see About the Cover.}
\expandafter\def\csname occovercredit@15\endcsname{Cover illustration: pair correlation of the zeta zeros against the Gaussian unitary ensemble; see About the Cover.}
\newcommand{\covercredit}{\ifcsname occovercredit@\VOLUMENUM\endcsname\csname occovercredit@\VOLUMENUM\endcsname\par\fi}

\newcommand{\bookauthors}{SEGHAIER Oussama}
\newcommand{\bookcopyright}{Copyright \textcopyright\ 2026 SEGHAIER Oussama}
\newcommand{\booklicence}{All rights reserved. No part of this book may be reproduced, stored
  in a retrieval system or transmitted in any form or by any means without the prior written
  permission of the copyright holder, except for brief quotations in reviews and scholarly works.}
% The imprint KDP records for a book with a free KDP ISBN.
\newcommand{\bookpublisher}{Independently published}
% Per-volume settings: a volume may \renewcommand these in its main.tex, after \usepackage[T1]{fontenc}
\usepackage[utf8]{inputenc}
\usepackage{mathpazo}
% Typewriter text in the txfonts typewriter, a vector font (the default EC one is a bitmap without cm-super),
% so that every font is an embedded Type 1 font, as print shops ask.
\renewcommand{\ttdefault}{txtt}
% Embedded vector sans-serif labels, consistent with the covers.
\renewcommand{\sfdefault}{qhv}
\usepackage{amsmath,amssymb,amsthm,mathtools}
% Every PDF uses the standard 7 x 10 inch book layout. Its identical file also
% serves as the paperback interior; the full cover is a separate publishing asset.
% Generated by tools/print_layout.py; edit print-layout.json.
\newcommand{\PRINTWIDTH}{7}
\newcommand{\PRINTHEIGHT}{10}
\newcommand{\PRINTBLEED}{0.125}
\newcommand{\PRINTSPINECOEFF}{0.002347}

\providecommand{\PRINTINNER}{0.75in}
% Keep these hooks for volume layouts and stacked mathematical figures.
\newif\ifprintformat\printformattrue
\newif\ifsmallformat\smallformattrue
\makeatletter\@twosidetrue\@mparswitchtrue\@openrighttrue\makeatother
\usepackage[paperwidth=\PRINTWIDTH in,paperheight=\PRINTHEIGHT in,twoside,inner=\PRINTINNER,outer=0.6in,
  top=0.75in,bottom=0.8in,headsep=0.2in,footskip=0.4in]{geometry}
% Break long formula chains for the book's text measure and constrain displays.
\newcommand{\qq}{\\}\newcommand{\eqbr}{\\&}
\AtBeginDocument{%
  \def\[#1\]{\begin{equation*}\sbox0{$\displaystyle#1$}%
    \ifdim\wd0>\linewidth\typeout{SHRINK \the\wd0\space > \the\linewidth\space at line \the\inputlineno}%
      \resizebox{\linewidth}{!}{\usebox0}\else\usebox0\fi\end{equation*}}}
\usepackage[final]{microtype}
\setlength{\emergencystretch}{3em}
\usepackage{enumitem}
\usepackage[dvipsnames]{xcolor}
% Approved charcoal/coral editorial identity. See docs/DESIGN.md.
\definecolor{ocBlack}{HTML}{191919}
\definecolor{ocCoral}{HTML}{F28B79}
\definecolor{ocCoralDeep}{HTML}{A34232}
\definecolor{ocInk}{HTML}{252323}
\definecolor{ocCream}{HTML}{F5F1E9}
% Compatibility names for mathematical figures authored before the palette change.
\colorlet{ocGold}{ocCoral}
\colorlet{ocGoldDeep}{ocCoralDeep}
% Legacy names remain aliases so semantic boxes and existing figures stay compatible.
\colorlet{ocTerracotta}{ocCoralDeep}
\colorlet{ocRust}{ocCoralDeep}
\definecolor{ocSepia}{HTML}{795B54}
\definecolor{ocOlive}{HTML}{5D6D65}
\definecolor{ocUmber}{HTML}{746072}
\colorlet{ocExercise}{ocSepia}
\colorlet{ocReminder}{ocCoralDeep}
\colorlet{ocBeauty}{ocOlive}
\colorlet{ocThink}{ocUmber}

% Paperback interiors are black ink on white paper; tint hierarchy survives in gray.
\ifprintformat
  \colorlet{ocExercise}{black!65}
  \colorlet{ocReminder}{black!70}
  \colorlet{ocBeauty}{black!60}
  \colorlet{ocThink}{black!75}
\fi
\usepackage[most]{tcolorbox}
\usepackage{needspace}
\usepackage{currfile}
% Book pages keep framed material whole. Record natural heights for the print
% audit; oversized material must be reorganised in the source, never clipped.
\makeatletter
\newwrite\ocLayoutFile
\AtBeginDocument{\immediate\openout\ocLayoutFile=\jobname.oclayout\relax}
\AtEndDocument{\immediate\closeout\ocLayoutFile}
\newcommand{\ocCheckBox}{%
  \immediate\write\ocLayoutFile{BOX|\currfilepath|\the\inputlineno|\tcb@height|\the\textheight}%
  \ifdim\tcb@height>\textheight
    \PackageError{orderchance}{Box exceeds the 7x10 text height at \currfilepath:\the\inputlineno}%
      {Reorganise this material into shorter, logically complete boxes. Do not shrink the type.}%
  \fi}
\newcommand{\ocSplitBoxError}{%
  \PackageError{orderchance}{A framed box was split across pages}%
    {Use an unbreakable box and reorganise oversized source material.}}
\tcbset{finish unbroken={\ocCheckBox},
  finish first={\ocSplitBoxError},finish middle={\ocSplitBoxError},finish last={\ocSplitBoxError}}
% Needspace inserts a break opportunity. Preserve LaTeX's heading keep when the
% first hint or solution immediately follows a section heading.
\newcommand{\ocOpeningSpace}{\if@nobreak\else\Needspace{3\baselineskip}\fi}
\makeatother
% Avoid isolated opening/closing lines and an isolated line before a display.
\clubpenalty=10000
\widowpenalty=10000
\displaywidowpenalty=10000
\raggedbottom
% Packages and macros used by figures; shared by the books and tools/figure.tex.
\usepackage{tikz}
\usetikzlibrary{arrows.meta,calc,positioning,decorations.pathreplacing}
\usepackage{pgfplots}
\pgfplotsset{compat=1.18}

\usepackage{pdfpages}
\usepackage{adjustbox}
% The print interior starts with its half-title. The build attaches the approved
% front to a separate reading PDF; EPUB and paperback wraps have their own covers.
% Retain the manuscript hook so every main.tex keeps the same front-matter order.
\newcommand{\makecover}{}
% Running heads in print: the chapter title on left pages, the section on right pages, in small
% italics, cut short with an ellipsis when too long for the text measure.
\ifprintformat
  \usepackage{fancyhdr}
  \usepackage[fit]{truncate}
  \pagestyle{fancy}\fancyhf{}
  \renewcommand{\headrulewidth}{0pt}
  \renewcommand{\chaptermark}[1]{\markboth{#1}{#1}}
  \renewcommand{\sectionmark}[1]{\markright{\thesection\quad #1}}
  \fancyhead[LE]{\small\thepage\hfill\truncate{0.85\headwidth}{\itshape\leftmark}}
  \fancyhead[RO]{\small\truncate{0.85\headwidth}{\itshape\rightmark}\hfill\thepage}
  % The blank page before a chapter has no running head.
  \makeatletter
  \renewcommand{\cleardoublepage}{\clearpage\if@twoside\ifodd\c@page\else
    \hbox{}\thispagestyle{empty}\newpage\fi\fi}
  \makeatother
\fi
% PDF links are black; online and EPUB themes are configured separately.
\ifprintformat
  \usepackage[hidelinks]{hyperref}
\else
  \usepackage[colorlinks=true,linkcolor=ocRust,urlcolor=ocRust]{hyperref}
\fi
% Half-title, series page, title page and copyright page (\maketitle builds all four).
% Front matter of every volume, in the order of a published book. \maketitle (called by each
% main.tex) builds four pages: the half-title (recto), the series page (verso), the title page
% (recto) and the copyright page (verso). The preface and the table of contents follow.
%
% Series-wide settings (one place for every volume):
\newcommand{\seriesname}{Order and Chance}
\newcommand{\seriessubtitle}{The Hidden Architecture of Mathematics}
\input{../common/volume-subtitles}
\input{../common/cover-credits}
\newcommand{\bookauthors}{SEGHAIER Oussama}
\newcommand{\bookcopyright}{Copyright \textcopyright\ 2026 SEGHAIER Oussama}
\newcommand{\booklicence}{All rights reserved. No part of this book may be reproduced, stored
  in a retrieval system or transmitted in any form or by any means without the prior written
  permission of the copyright holder, except for brief quotations in reviews and scholarly works.}
% The imprint KDP records for a book with a free KDP ISBN.
\newcommand{\bookpublisher}{Independently published}
% Per-volume settings: a volume may \renewcommand these in its main.tex, after \input{../common/preamble}.
%   \bookedition    the edition line
%   \bookisbnprint  the paperback ISBN (KDP assigns one free of charge)
%   \bookisbnebook  the eBook ISBN (optional on Kindle, which uses an ASIN)
\newcommand{\bookedition}{First edition, 2026}
\newcommand{\bookisbnprint}{to be assigned}
\newcommand{\bookisbnebook}{}
\makeatletter\newcommand{\oc@empty}{}\makeatother
% The PDF's author metadata, whatever a volume's main.tex sets.
\AtBeginDocument{\hypersetup{pdfauthor={SEGHAIER Oussama}}}
% Volume number and title, passed by build.sh from the "% epub-title:" line of main.tex.
\providecommand{\VOLUMENUM}{}
\providecommand{\VOLUMETITLE}{}

\makeatletter
% The series list, read from common/series.tex; the current volume is set in semibold.
\newcommand{\seriespart}[1]{\par\medskip\noindent{\ocLabelFont\ocLab{#1}}\par\nobreak\smallskip}
\newcommand{\seriesvolume}[3]{%
  \def\oc@this{#1}%
  \par\noindent\hangindent=2.2em\makebox[2.2em][l]{\small\ocNum{#1}}%
  \ifx\oc@this\VOLUMENUM{\small\textbf{#2}}\else{\small #2}\fi
  \if\relax\detokenize{#3}\relax\else{\footnotesize\itshape\ --- #3}\fi\par}
% Front pages in the interior's system: Cormorant titles, Fira labels, EB Garamond text, a
% coral rule. Text is set flush left on the text block, like the chapter openers.
\newcommand{\oc@masthead}[1]{{\ocDisplay\fontsize{#1}{#1}\selectfont\color{ocBlack}\ocCaps[26]{\seriesname}}}
\renewcommand{\maketitle}{%
  \cleardoublepage
  % Half-title
  \thispagestyle{empty}%
  \vspace*{0.22\textheight}%
  {\raggedright
   \oc@masthead{12}\par\vspace{22pt}%
   \ifx\VOLUMETITLE\@empty\else
     {\ocLabelFont\ocLab{Volume}\enspace\ocNum{\VOLUMENUM}}\par\vspace{8pt}%
     {\ocDisplay\linespread{1}\fontsize{26}{30}\selectfont\color{ocBlack}\VOLUMETITLE\par}\vspace{12pt}%
   \fi
   {\color{ocCoral}\rule{0.9in}{1.2pt}}\par}%
  \clearpage
  % Series page
  \thispagestyle{empty}%
  {\raggedright\let\oc@this\relax
   \noindent\oc@masthead{11}\par\vspace{4pt}%
   \noindent{\small\itshape\seriessubtitle}\par\medskip
   \input{../common/series}\par}%
  \clearpage
  % Title page
  \thispagestyle{empty}%
  \vspace*{0.12\textheight}%
  {\raggedright
   \oc@masthead{12}\par\vspace{0.09\textheight}%
   \ifx\VOLUMETITLE\@empty
     {\ocDisplay\linespread{1}\fontsize{30}{36}\selectfont\color{ocBlack}\let\textbf\relax\let\Huge\relax
      \let\LARGE\relax\let\Large\relax\@title\par}%
   \else
     {\ocLabelFont\ocLab{Volume}\enspace\ocNum{\VOLUMENUM}}\par\vspace{10pt}%
     {\ocDisplay\linespread{1}\fontsize{38}{42}\selectfont\color{ocBlack}\hyphenpenalty10000 \VOLUMETITLE\par}%
     \vspace{12pt}{\large\itshape\volumesubtitle{\VOLUMENUM}\par}%
   \fi
   \vspace{16pt}{\color{ocCoral}\rule{0.9in}{1.2pt}}\par\vspace{0.06\textheight}%
   {\ocDisplay\fontsize{15}{18}\selectfont\color{ocBlack}\ocCaps[12]{\bookauthors}\par}%
   \vfill
   {\ocBook\fontsize{8.5}{11}\selectfont\color{ocBlack}\bookedition\par}}%
  \clearpage
  % Copyright page: the text sits at the foot of the page, as is the custom.
  \thispagestyle{empty}%
  \vspace*{\fill}%
  {\footnotesize\raggedright\setlength{\parskip}{0.6em}\setlength{\parindent}{0pt}%
   \textit{\seriesname}\ifx\VOLUMETITLE\@empty\else. Volume~\VOLUMENUM: \textit{\VOLUMETITLE}\fi\par
   \bookcopyright\par
   \booklicence\par
   \bookedition. \bookpublisher. This printing typeset on \today.\par
   ISBN (paperback): \bookisbnprint\ifx\bookisbnebook\oc@empty\else\\ISBN (eBook): \bookisbnebook\fi\par
   \covercredit
   Typeset with XeLaTeX in EB Garamond and Garamond-Math, with Cormorant Garamond headings
   and Fira Sans labels.\par}%
  \clearpage
  \global\let\maketitle\relax}
\makeatother


\newtheorem{theorem}{Theorem}[chapter]
\newtheorem{proposition}[theorem]{Proposition}
\newtheorem{lemma}[theorem]{Lemma}
\newtheorem{corollary}[theorem]{Corollary}
\theoremstyle{definition}
\newtheorem{definition}[theorem]{Definition}
\newtheorem{example}[theorem]{Example}
\theoremstyle{remark}
\newtheorem*{remark}{Remark}

% Exercises and solutions
\newcounter{exo}[chapter]
\renewcommand{\theexo}{\thechapter.\arabic{exo}}
% The counter anchor travels inside the box when it moves to the next page.
\tcbset{oc exercise/.style={code={\stepcounter{exo}},
  before upper={\addtocounter{exo}{-1}\refstepcounter{exo}}}}
\newtcolorbox{exercise}[1][]{%
    oc exercise,enhanced,unbreakable,colback=ocExercise!6,colframe=ocExercise,
    boxrule=0.6pt,arc=2pt,left=6pt,right=6pt,
    title={\textbf{Exercise \theexo}\ifx\relax#1\relax\else\ \textit{(#1)}\fi}}
\newenvironment{solution}[1]{%
  \par\ocOpeningSpace\medskip\noindent\textbf{\color{ocReminder}Solution of Exercise~\ref{#1}.}\quad}%
  {\hfill$\blacksquare$\par\medskip}
\newtcolorbox{reminder}[1][]{enhanced,unbreakable,colback=ocReminder!4,colframe=ocReminder,
  boxrule=0.6pt,arc=2pt,title={\textbf{Reminder}\ifx\relax#1\relax\else: #1\fi}}
\newtcolorbox{beauty}{enhanced,unbreakable,colback=ocBeauty!5,colframe=ocBeauty,
  boxrule=0.6pt,arc=2pt,title={\textbf{Why this is beautiful}}}

% Pause and think: a question to try before reading on.
\newtcolorbox{think}{enhanced,unbreakable,colback=ocThink!5,colframe=ocThink,
  boxrule=0.6pt,arc=2pt,title={\textbf{Pause and think}}}
% Hints, printed between the exercises and the solutions.
\newenvironment{hint}[1]{%
  \par\ocOpeningSpace\smallskip\noindent\textbf{\color{ocThink}Hint for Exercise~\ref{#1}.}\quad}%
  {\par\smallskip}
% Figures: \fig{name} draws figures/name.tex (a tikzpicture); the EPUB uses an SVG of it.
\newcommand{\fig}[1]{\adjustbox{max width=\linewidth}{\input{figures/#1}}}

\newcommand{\N}{\mathbb{N}}
\newcommand{\Z}{\mathbb{Z}}
\newcommand{\Q}{\mathbb{Q}}
\newcommand{\R}{\mathbb{R}}
\newcommand{\C}{\mathbb{C}}
\newcommand{\dd}{\,\mathrm{d}}
% Stirling numbers (Tome 4): \stirlingii nk (partitions), \stirlingi nk (cycles).
\newcommand{\stirlingii}[2]{\genfrac{\{}{\}}{0pt}{}{#1}{#2}}
\newcommand{\stirlingi}[2]{\genfrac{[}{]}{0pt}{}{#1}{#2}}
\setlist{itemsep=2pt}
% Solutions: \idea explains how to find the approach, \writeup gives the full argument.
\newcommand{\idea}{\par\nobreak\noindent\textit{How to find it.}\quad}
\newcommand{\writeup}{\par\smallskip\noindent\textit{Write-up.}\quad}
.
%   \bookedition    the edition line
%   \bookisbnprint  the paperback ISBN (KDP assigns one free of charge)
%   \bookisbnebook  the eBook ISBN (optional on Kindle, which uses an ASIN)
\newcommand{\bookedition}{First edition, 2026}
\newcommand{\bookisbnprint}{to be assigned}
\newcommand{\bookisbnebook}{}
\makeatletter\newcommand{\oc@empty}{}\makeatother
% The PDF's author metadata, whatever a volume's main.tex sets.
\AtBeginDocument{\hypersetup{pdfauthor={SEGHAIER Oussama}}}
% Volume number and title, passed by build.sh from the "% epub-title:" line of main.tex.
\providecommand{\VOLUMENUM}{}
\providecommand{\VOLUMETITLE}{}

\makeatletter
% The series list, read from common/series.tex; the current volume is set in semibold.
\newcommand{\seriespart}[1]{\par\medskip\noindent{\ocLabelFont\ocLab{#1}}\par\nobreak\smallskip}
\newcommand{\seriesvolume}[3]{%
  \def\oc@this{#1}%
  \par\noindent\hangindent=2.2em\makebox[2.2em][l]{\small\ocNum{#1}}%
  \ifx\oc@this\VOLUMENUM{\small\textbf{#2}}\else{\small #2}\fi
  \if\relax\detokenize{#3}\relax\else{\footnotesize\itshape\ --- #3}\fi\par}
% Front pages in the interior's system: Cormorant titles, Fira labels, EB Garamond text, a
% coral rule. Text is set flush left on the text block, like the chapter openers.
\newcommand{\oc@masthead}[1]{{\ocDisplay\fontsize{#1}{#1}\selectfont\color{ocBlack}\ocCaps[26]{\seriesname}}}
\renewcommand{\maketitle}{%
  \cleardoublepage
  % Half-title
  \thispagestyle{empty}%
  \vspace*{0.22\textheight}%
  {\raggedright
   \oc@masthead{12}\par\vspace{22pt}%
   \ifx\VOLUMETITLE\@empty\else
     {\ocLabelFont\ocLab{Volume}\enspace\ocNum{\VOLUMENUM}}\par\vspace{8pt}%
     {\ocDisplay\linespread{1}\fontsize{26}{30}\selectfont\color{ocBlack}\VOLUMETITLE\par}\vspace{12pt}%
   \fi
   {\color{ocCoral}\rule{0.9in}{1.2pt}}\par}%
  \clearpage
  % Series page
  \thispagestyle{empty}%
  {\raggedright\let\oc@this\relax
   \noindent\oc@masthead{11}\par\vspace{4pt}%
   \noindent{\small\itshape\seriessubtitle}\par\medskip
   % Generated from series-catalog.json by tools/series_metadata.py.
% The series list for every series page and back cover.
\seriespart{Part I. Foundations}
\seriesvolume{1}{First Principles}{Proof, Problem Solving and Mathematical Thinking}
\seriesvolume{2}{The Integers}{Primes, Patterns and Proofs}
\seriesvolume{3}{Symmetry}{Groups, Rings and Fields}
\seriesvolume{4}{Chance}{The Logic of Uncertainty}
\seriespart{Part II. Bridges}
\seriesvolume{5}{Calculation}{From Formula to Insight}
\seriesvolume{6}{Crossroads}{Arithmetic Meets Analysis}
\seriespart{Part III. Analysis and Geometry}
\seriesvolume{7}{Spaces}{Topology, Measure and Limits}
\seriesvolume{8}{Analysis}{Operators, Spectra and Transforms}
\seriesvolume{9}{Shapes}{The Geometry of Connection}
\seriesvolume{10}{Abstraction}{Categories, Algebra and Homology}
\seriespart{Part IV. Randomness}
\seriesvolume{11}{Randomness}{Martingales, Limits and Ergodicity}
\seriesvolume{12}{Random Motion}{The Geometry of Random Paths}
\seriespart{Part V. Number}
\seriesvolume{13}{Arithmetic}{From Equations to Hidden Worlds}
\seriesvolume{14}{The Primes}{Their Hidden Rhythm}
\seriespart{Part VI. Synthesis}
\seriesvolume{15}{Order and Chance}{The Meeting of Structure and Randomness}
\par}%
  \clearpage
  % Title page
  \thispagestyle{empty}%
  \vspace*{0.12\textheight}%
  {\raggedright
   \oc@masthead{12}\par\vspace{0.09\textheight}%
   \ifx\VOLUMETITLE\@empty
     {\ocDisplay\linespread{1}\fontsize{30}{36}\selectfont\color{ocBlack}\let\textbf\relax\let\Huge\relax
      \let\LARGE\relax\let\Large\relax\@title\par}%
   \else
     {\ocLabelFont\ocLab{Volume}\enspace\ocNum{\VOLUMENUM}}\par\vspace{10pt}%
     {\ocDisplay\linespread{1}\fontsize{38}{42}\selectfont\color{ocBlack}\hyphenpenalty10000 \VOLUMETITLE\par}%
     \vspace{12pt}{\large\itshape\volumesubtitle{\VOLUMENUM}\par}%
   \fi
   \vspace{16pt}{\color{ocCoral}\rule{0.9in}{1.2pt}}\par\vspace{0.06\textheight}%
   {\ocDisplay\fontsize{15}{18}\selectfont\color{ocBlack}\ocCaps[12]{\bookauthors}\par}%
   \vfill
   {\ocBook\fontsize{8.5}{11}\selectfont\color{ocBlack}\bookedition\par}}%
  \clearpage
  % Copyright page: the text sits at the foot of the page, as is the custom.
  \thispagestyle{empty}%
  \vspace*{\fill}%
  {\footnotesize\raggedright\setlength{\parskip}{0.6em}\setlength{\parindent}{0pt}%
   \textit{\seriesname}\ifx\VOLUMETITLE\@empty\else. Volume~\VOLUMENUM: \textit{\VOLUMETITLE}\fi\par
   \bookcopyright\par
   \booklicence\par
   \bookedition. \bookpublisher. This printing typeset on \today.\par
   ISBN (paperback): \bookisbnprint\ifx\bookisbnebook\oc@empty\else\\ISBN (eBook): \bookisbnebook\fi\par
   \covercredit
   Typeset with XeLaTeX in EB Garamond and Garamond-Math, with Cormorant Garamond headings
   and Fira Sans labels.\par}%
  \clearpage
  \global\let\maketitle\relax}
\makeatother


\newtheorem{theorem}{Theorem}[chapter]
\newtheorem{proposition}[theorem]{Proposition}
\newtheorem{lemma}[theorem]{Lemma}
\newtheorem{corollary}[theorem]{Corollary}
\theoremstyle{definition}
\newtheorem{definition}[theorem]{Definition}
\newtheorem{example}[theorem]{Example}
\theoremstyle{remark}
\newtheorem*{remark}{Remark}

% Exercises and solutions
\newcounter{exo}[chapter]
\renewcommand{\theexo}{\thechapter.\arabic{exo}}
% The counter anchor travels inside the box when it moves to the next page.
\tcbset{oc exercise/.style={code={\stepcounter{exo}},
  before upper={\addtocounter{exo}{-1}\refstepcounter{exo}}}}
\newtcolorbox{exercise}[1][]{%
    oc exercise,enhanced,unbreakable,colback=ocExercise!6,colframe=ocExercise,
    boxrule=0.6pt,arc=2pt,left=6pt,right=6pt,
    title={\textbf{Exercise \theexo}\ifx\relax#1\relax\else\ \textit{(#1)}\fi}}
\newenvironment{solution}[1]{%
  \par\ocOpeningSpace\medskip\noindent\textbf{\color{ocReminder}Solution of Exercise~\ref{#1}.}\quad}%
  {\hfill$\blacksquare$\par\medskip}
\newtcolorbox{reminder}[1][]{enhanced,unbreakable,colback=ocReminder!4,colframe=ocReminder,
  boxrule=0.6pt,arc=2pt,title={\textbf{Reminder}\ifx\relax#1\relax\else: #1\fi}}
\newtcolorbox{beauty}{enhanced,unbreakable,colback=ocBeauty!5,colframe=ocBeauty,
  boxrule=0.6pt,arc=2pt,title={\textbf{Why this is beautiful}}}

% Pause and think: a question to try before reading on.
\newtcolorbox{think}{enhanced,unbreakable,colback=ocThink!5,colframe=ocThink,
  boxrule=0.6pt,arc=2pt,title={\textbf{Pause and think}}}
% Hints, printed between the exercises and the solutions.
\newenvironment{hint}[1]{%
  \par\ocOpeningSpace\smallskip\noindent\textbf{\color{ocThink}Hint for Exercise~\ref{#1}.}\quad}%
  {\par\smallskip}
% Figures: \fig{name} draws figures/name.tex (a tikzpicture); the EPUB uses an SVG of it.
\newcommand{\fig}[1]{\adjustbox{max width=\linewidth}{\input{figures/#1}}}

\newcommand{\N}{\mathbb{N}}
\newcommand{\Z}{\mathbb{Z}}
\newcommand{\Q}{\mathbb{Q}}
\newcommand{\R}{\mathbb{R}}
\newcommand{\C}{\mathbb{C}}
\newcommand{\dd}{\,\mathrm{d}}
% Stirling numbers (Tome 4): \stirlingii nk (partitions), \stirlingi nk (cycles).
\newcommand{\stirlingii}[2]{\genfrac{\{}{\}}{0pt}{}{#1}{#2}}
\newcommand{\stirlingi}[2]{\genfrac{[}{]}{0pt}{}{#1}{#2}}
\setlist{itemsep=2pt}
% Solutions: \idea explains how to find the approach, \writeup gives the full argument.
\newcommand{\idea}{\par\nobreak\noindent\textit{How to find it.}\quad}
\newcommand{\writeup}{\par\smallskip\noindent\textit{Write-up.}\quad}
